{\small
\begin{longtable}{>{\centering\arraybackslash}p{0.070\textwidth}>{\centering\arraybackslash}p{0.150\textwidth}>{\centering\arraybackslash}p{0.240\textwidth}>{\centering\arraybackslash}p{0.160\textwidth}>{\centering\arraybackslash}p{0.360\textwidth}}
  \caption{变量编码与变换方案}
  \label{tab:encode} \\
  \toprule
  变量 & 识别类型 & 编码/变换 & 基线/参照 & 解释 \\
  \midrule
  \endfirsthead
  \caption{变量编码与变换方案（续）} \\
  \toprule
  变量 & 识别类型 & 编码/变换 & 基线/参照 & 解释 \\
  \midrule
  \endhead
  \bottomrule
  \endfoot
  X1 & 连续 & log(X1) & — & 粉丝规模弹性；避免量级差异 \\
  X2 & 名义（4 类） & 独热编码，3 个哑变量 & Photo & 比 Photo 的相对曝光差异 \\
  X3 & 名义（3 类） & 独热编码，2 个哑变量 & 1 & 比主题 1 的相对曝光差异 \\
  X4 & 有序/周期（12） & 独热编码，11 个哑变量 & 1 月 & 控制季节性与时间趋势 \\
  X5 & 有序/周期（7） & 独热编码，6 个哑变量 & 周一 & 控制星期效应 \\
  X6 & 有序/周期（24） & 按分布合并为 3 组，2 个哑变量 & 上午（6–11 时） & 凌晨 0–5、上午 6–11、下午晚间 12–23；避免稀疏单元 \\
  X7 & 二元 & 0/1 数值 & 0（自然） & 付费的额外曝光 \\
  Y1 & 连续计数 & log1p(Y1) & — & 压缩重尾，系数近似百分比效应 \\
\end{longtable}
}
